\documentclass[aspectratio=169]{beamer}
\input{header}
\coursecode{JKSSB}

\title{Current Office, Email \& Statement Evaluation}
\subtitle{Lesson 61 --- Advanced Statements and Near-Neighbour Traps}
\author{JKSSB Computer Awareness}
\date{\today}

\begin{document}
\frame{\titlepage}

\begin{frame}{Word: Structured Content and Content Controls}
  \begin{itemize}
    \item A \key{content control} is a building block inserted into a
      document or template to build a structured form.
    \item Common types: plain text, rich text, picture, drop-down list, combo
      box, date picker, and check box.
    \item They are inserted from the \key{Developer} tab (enable it in Word
      Options $\rightarrow$ Customize Ribbon).
    \item A content control can \key{restrict editing} and enforce limited
      \key{data validation} (for example a drop-down list or a date picker).
  \end{itemize}
  \begin{block}{One line}
    Content controls let a template allow only certain kinds of entry at a
    chosen spot in the document.
  \end{block}
\end{frame}

\begin{frame}{First-Line Indent vs Hanging Indent}
  \begin{center}
    \begin{tabular}{p{0.20\textwidth}p{0.36\textwidth}p{0.34\textwidth}}
      \toprule
      \textbf{Indent} & \textbf{What it shifts} & \textbf{Typical use} \\
      \midrule
      First-line indent & Only the \key{first line} of the paragraph & Normal
      paragraph openings \\
      Hanging indent & Every line \key{except the first} (the first line hangs
      to the left) & References, bibliography, bulleted lists \\
      \bottomrule
    \end{tabular}
  \end{center}
  \vspace{0.25cm}
  \muted{These are paragraph-format settings; two paragraphs can use different
  indents even under the same style.}
\end{frame}

\begin{frame}{The Navigation Pane}
  \begin{itemize}
    \item Open it from the \key{View} tab $\rightarrow$ Navigation Pane.
    \item It offers three views: \key{Headings}, \key{Pages}, and
      \key{Results} (search results).
    \item It is used to \key{jump} quickly to a heading or a page and to
      \key{search} the document.
    \item The Headings list is built from paragraphs that carry a
      \key{Heading style} (Heading 1, Heading 2, \ldots).
  \end{itemize}
  \begin{alertblock}{Trap}
    Reviewers' \key{comments} are read in the markup/comments pane,
    \bad{not} directly in the Navigation Pane. (PYQ-12)
  \end{alertblock}
\end{frame}

\begin{frame}{Track Changes and Its Metadata}
  \begin{itemize}
    \item \key{Track Changes} records \key{insertions}, \key{deletions}, and
      \key{formatting} changes.
    \item Each recorded change stores the \key{author's name} and the
      \key{date and time}.
    \item \key{Accept} a change: it becomes final text and its revision
      metadata is \key{removed}.
    \item \key{Reject} a change: the original text is restored.
  \end{itemize}
  \begin{block}{Remember}
    Accepting a change keeps the new text but strips the author metadata
    attached to that change. (PYQ-12)
  \end{block}
\end{frame}

\begin{frame}{Section Break vs Page Break}
  \begin{center}
    \begin{tabular}{p{0.20\textwidth}p{0.34\textwidth}p{0.36\textwidth}}
      \toprule
      \textbf{Break} & \textbf{What it does} & \textbf{What changes} \\
      \midrule
      Page Break & Moves text to the next page & Nothing: same section, same
      headers and margins \\
      Section Break & Starts a \key{new section} & Each section can have its
      own headers/footers, margins, orientation, columns, numbering \\
      \bottomrule
    \end{tabular}
  \end{center}
  \vspace{0.25cm}
  \begin{alertblock}{Trap}
    Different headers and footers in one document require a \key{Section
    Break}, \bad{not} a Page Break. (PYQ-13, TRAP-13)
  \end{alertblock}
\end{frame}

\begin{frame}{Excel: The Nested IF}
  \begin{itemize}
    \item An \key{IF} returns one of two values:
      \texttt{=IF(condition, value\_if\_true, value\_if\_false)}.
    \item A \key{nested IF} places another IF in one of those two positions.
    \item \texttt{=IF(A1>50,"High",IF(A1>25,"Medium","Low"))}
    \item The \key{first} condition that is TRUE decides the result; testing
      stops there.
  \end{itemize}
  \begin{exampleblock}{Worked example}
    A1\,$=$\,70 $\rightarrow$ \good{High}; A1\,$=$\,30 $\rightarrow$
    \good{Medium}; A1\,$=$\,10 $\rightarrow$ \good{Low}.
  \end{exampleblock}
\end{frame}

\begin{frame}{Blank and Non-Numeric Input in IF}
  \begin{itemize}
    \item In a numeric comparison, a \key{blank cell is treated as zero}.
    \item So \texttt{=IF(A1>50,\ldots)} with A1 blank takes the \key{FALSE}
      branch, because 0 is not greater than 50.
    \item Non-numeric \key{text} is not a number: arithmetic on it gives
      \texttt{\#VALUE!}, and \texttt{ISNUMBER} returns FALSE.
    \item Guard the formula with \key{ISBLANK}, \key{ISNUMBER}, or
      \key{IFERROR} when the cell may be empty or text.
  \end{itemize}
  \begin{alertblock}{Trap}
    A blank cell does \bad{not} make IF stop or error --- it is read as 0.
    (PYQ-14)
  \end{alertblock}
\end{frame}

\begin{frame}{COUNTIFS: Multi-Criterion Counting}
  \begin{itemize}
    \item \key{COUNTIFS} counts the rows that satisfy \key{all} the given
      conditions (AND logic).
    \item Form: \texttt{COUNTIFS(range1, criteria1, range2, criteria2,
      \ldots)}.
    \item All ranges must be the \key{same size}.
    \item Use \key{COUNTIF} (singular) for one condition; \key{COUNTIFS}
      (plural) for several.
    \item Text criteria may use the wildcards \texttt{*} and \texttt{?}.
  \end{itemize}
  \begin{block}{Example}
    \texttt{=COUNTIFS(A2:A20,"Sales",B2:B20,">1000")} counts the Sales rows
    above 1000.
  \end{block}
\end{frame}

\begin{frame}{SUMPRODUCT: Products and Counting}
  \begin{itemize}
    \item \key{SUMPRODUCT} multiplies corresponding elements of the arrays and
      returns the \key{sum of the products}.
    \item \texttt{=SUMPRODUCT(A1:A3,B1:B3)} =
      \texttt{A1*B1 + A2*B2 + A3*B3}.
    \item With boolean arrays it can also \key{count}:
      \texttt{=SUMPRODUCT((A1:A10>50)*(B1:B10="Yes"))} counts rows meeting
      both conditions.
  \end{itemize}
  \begin{alertblock}{Trap}
    \key{COUNTIFS} takes several criteria; \key{SUMPRODUCT} does \bad{not}
    ``only add'' --- it can count multi-criterion rows as well. (PYQ-15,
    TRAP-14)
  \end{alertblock}
\end{frame}

\begin{frame}{Email Filters and Rules}
  \begin{itemize}
    \item A \key{filter} (called a \key{rule} in Outlook) acts
      \key{automatically} on incoming or outgoing mail.
    \item It checks \key{conditions}: sender, recipient, subject, keywords,
      attachments.
    \item It performs an \key{action}: move to a folder, flag, forward,
      delete, or mark as read.
    \item Typical use: route mail from one sender into a dedicated folder.
  \end{itemize}
  \begin{block}{One line}
    A rule = a condition plus an action, applied automatically to mail.
  \end{block}
\end{frame}

\begin{frame}{Digital Signature vs Encryption}
  \begin{center}
    \begin{tabular}{p{0.22\textwidth}p{0.36\textwidth}p{0.32\textwidth}}
      \toprule
      \textbf{Point} & \textbf{Digital signature} & \textbf{Encryption} \\
      \midrule
      Purpose & Authenticity and integrity & Confidentiality \\
      What it proves & Who sent it, and that it was not altered & Only the
      intended reader can read it \\
      Hides the content? & \bad{No} & \good{Yes} \\
      \bottomrule
    \end{tabular}
  \end{center}
  \vspace{0.25cm}
  \begin{alertblock}{Trap}
    A digital signature does \bad{not} encrypt the message body. It verifies
    the sender and detects tampering. (PYQ-16, TRAP-15)
  \end{alertblock}
\end{frame}

\begin{frame}{POP3: Post Office Protocol version 3}
  \begin{itemize}
    \item \key{POP3} stands for \key{Post Office Protocol version 3}.
    \item It \key{downloads} mail from the mail server to the \key{local}
      device.
    \item By default it typically \key{removes} the mail from the server, so
      it is not kept in sync across devices.
    \item \key{IMAP} keeps mail on the server and syncs all devices.
  \end{itemize}
  \begin{alertblock}{Trap}
    \key{SMTP} sends mail; \key{POP3} receives (downloads) it. POP3 does
    \bad{not} keep mail synced on the server. (PYQ-16, TRAP-17)
  \end{alertblock}
\end{frame}

\begin{frame}{BCC and Recipient Privacy}
  \begin{itemize}
    \item \key{To}: the primary recipients --- visible to everyone.
    \item \key{Cc} (Carbon Copy): copied recipients --- visible to everyone.
    \item \key{Bcc} (Blind Carbon Copy): the addresses are \key{hidden} from
      the other recipients.
    \item The \key{sender} always sees all recipients; a Bcc recipient cannot
      see the other recipients.
  \end{itemize}
  \begin{alertblock}{Trap}
    \key{Bcc} hides recipients from \key{each other}, \bad{not} from the
    sender. (PYQ-16, TRAP-16)
  \end{alertblock}
\end{frame}

\begin{frame}{PowerPoint SmartArt Conversion}
  \begin{itemize}
    \item Convert a \key{bulleted list} of text into SmartArt: \key{Home}
      tab $\rightarrow$ Convert to SmartArt Graphic.
    \item Convert SmartArt \key{back to text}: SmartArt Tools Design
      $\rightarrow$ Convert $\rightarrow$ Convert to Text (you get a bulleted
      list; shapes and pictures are lost).
    \item \key{Convert to Shapes} breaks the SmartArt into ordinary shapes;
      after that it \key{cannot} be turned back into SmartArt.
  \end{itemize}
  \begin{block}{Remember}
    The direct source of ``Convert to SmartArt'' is \key{text}, so a picture
    or a chart is not converted to SmartArt in one step. (PYQ-17)
  \end{block}
\end{frame}

\begin{frame}{Trap Alert: Near-Neighbour Distinctions}
  \begin{itemize}
    \item Headers/footers $\rightarrow$ \key{Section Break}, not Page Break.
    \item Hanging indent shifts \key{every line except the first}.
    \item Blank cell in a comparison = \key{0}, not an error.
    \item \key{COUNTIFS} = several criteria; \key{SUMPRODUCT} can add and
      count.
    \item \key{Digital signature} $\neq$ encryption.
    \item \key{POP3} downloads; \key{IMAP} syncs.
    \item \key{Bcc} hides recipients from each other, not from the sender.
  \end{itemize}
\end{frame}

\begin{frame}{Lesson 61: Recall}
  \begin{itemize}
    \item A \key{content control} structures a template and can restrict
      entry; it is inserted from the \key{Developer} tab.
    \item A \key{first-line} indent shifts the first line; a \key{hanging}
      indent shifts every line except the first.
    \item The \key{Navigation Pane} jumps by headings, pages, or results;
      comments are \bad{not} shown there.
    \item \key{Track Changes} stores author and time; accepting a change
      removes that change's metadata.
    \item \key{Section Break} (not Page Break) allows different headers and
      footers.
    \item A nested \key{IF} returns the first TRUE branch; a \key{blank} cell
      is read as 0.
    \item \key{COUNTIFS} counts with AND logic; \key{SUMPRODUCT} multiplies
      and can count.
    \item \key{Digital signature} = authenticity; \key{encryption} =
      confidentiality; \key{POP3} downloads; \key{Bcc} hides recipients.
    \item SmartArt converts from text and back to text; shapes cannot revert
      to SmartArt.
  \end{itemize}
\end{frame}
\end{document}
